% !TeX root = ../main.tex

\chapter{Conclusioni e Sviluppi Futuri}
\label{cap:conclusioni}

Il progetto ha dimostrato la fattibilità di un'architettura logistica resiliente in cui l'automazione è affidata a un layer \textit{Agentic AI} mediato da MCP, con un orchestratore che funge da \textit{triage manager} per \texttt{HealthAgent} (piano infrastrutturale) e \texttt{LogisticAgent} (piano di business). La separazione tra LLM, layer MCP e infrastruttura Kubernetes ha permesso di soddisfare contemporaneamente osservabilità, sicurezza e scalabilità.

\textbf{Risultati principali.} (i) \textit{Reliability infrastrutturale}: graceful degradation, exponential backoff e rischedulazione K8s hanno garantito RTO contenuto (12\,s su InfluxDB) e rischedulazione immediata del broker MQTT senza interruzione del monitoraggio. (ii) \textit{Safety agentica}: \textit{Least Privilege}, contatori di iterazione, kill-switch backend e HITL hanno confinato l'autonomia dell'LLM; la soglia dei 6 droni mantiene il controllo umano sulle azioni ad alto impatto. (iii) \textit{Sostenibilità economica}: OPEX prevedibile in intervallo 160--650\,\$/mese, con scelta consapevole tra capacità di ragionamento e costo per ciclo.

\textbf{Limitazioni.} La latenza del ciclo agentico (\texttt{time.sleep(30)} + inferenza) non è adatta a load spike $<60$\,s; la regola peso/usura è oggi solo prompt-driven; la suite di chaos non misura \textit{loss rate} di telemetria né latenza end-to-end sotto stress; i valori economici sono basati su ipotesi di carico costante.

\textbf{Sviluppi futuri.} (i) Affiancare un HPA Kubernetes nativo come \emph{fast-path} per i picchi a breve orizzonte, lasciando all'agente le decisioni strategiche. (ii) Promuovere la regola peso/usura da \textit{soft policy} a vincolo deterministico nel layer MCP. (iii) Estendere la suite di chaos engineering con latency injection, misurazione della loss rate di InfluxDB e validazione del kill-switch in outage prolungati. (iv) Implementare un routing dinamico tra modelli LLM (Mistral Small per ordinario, Gemini 2.5 Pro per anomalie complesse) per massimizzare il ROA. (v) Validare il digital twin su dispositivi fisici o simulatori ad alta fedeltà.

In sintesi, il progetto mostra come un layer agentico mediato da MCP, accompagnato da guardrail di sicurezza e da una rigorosa strategia di osservabilità, possa elevare il livello di automazione di un'infrastruttura logistica critica senza sacrificare la prevedibilità operativa e il controllo umano sulle decisioni ad alto impatto.
